\heading{热学-模拟题2-期中}

\noteinfo{题目和答案由AI生成。}

\textbf{一、简答题}
\begin{question}
写出 Maxwell 速率分布下最概然速率 $v_{\rm mp}$。
\begin{solution}

\ansmath{v_{\rm mp}=\sqrt{\frac{2k_BT}{m}}.}
\end{solution}
\end{question}

\begin{question}
写出稀薄硬球气体的平均自由程 $\ell$。
\begin{solution}

若数密度为 $n$、碰撞截面为 $\sigma$，则
\ansmath{\ell=\frac{1}{\sqrt2\,n\sigma}.}
\end{solution}
\end{question}

\begin{question}
对具有 $f$ 个二次自由度的理想气体，写出 $C_V$ 与绝热指数 $\gamma$。
\begin{solution}

\ansmath{C_V=\frac{f}{2}R,
\qquad
\gamma=\frac{f+2}{f}.}
\end{solution}
\end{question}

\begin{question}
写出理想气体绝热声速的表达式。
\begin{solution}

\ansmath{c=\sqrt{\gamma\frac{p}{\rho}}.}
\end{solution}
\end{question}

\begin{question}
若气体处于温度为 $T$ 的平衡态，单个分子的平均平动动能是多少？
\begin{solution}

\ansmath{\langle \varepsilon\rangle=\frac32 k_BT.}
\end{solution}
\end{question}

\textbf{二、计算与证明题}
\begin{question}
一圆柱形刚性容器绕其轴线以角速度 $\omega$ 匀速旋转，内部为温度恒定的稀薄理想气体。设数密度只与到轴距离 $r$ 有关。
\begin{enumerate}[label=（\arabic*）,leftmargin=2.8em,itemsep=0.2em]
\item 求平衡数密度 $n(r)$；
\item 若容器半径为 $R$，求边壁与轴线上压强之比 $p(R)/p(0)$。
\end{enumerate}
\begin{solution}

\begin{enumerate}[label=（\arabic*）,leftmargin=2.8em,itemsep=0.2em]
\item 在旋转参考系中，单位质量受到离心势
\[
\Phi(r)=-\frac12\omega^2r^2.
\]
Boltzmann 分布给出
\[
n(r)=n(0)\exp\left(-\frac{m\Phi(r)}{k_BT}\right)
=n(0)\exp\left(\frac{m\omega^2r^2}{2k_BT}\right).
\]
故
\ansmath{n(r)=n(0)\exp\left(\frac{m\omega^2r^2}{2k_BT}\right).}

\item 由 $p=nk_BT$，有
\ansmath{\frac{p(R)}{p(0)}=
\exp\left(\frac{m\omega^2R^2}{2k_BT}\right).}
\end{enumerate}
\end{solution}
\end{question}

\begin{question}
由 Maxwell 速率分布推导单个分子的平动动能
\[
\varepsilon=\frac12 mv^2
\]
的概率密度函数，并求其均值与方差。
\begin{solution}

速率分布为
\[
f_v(v)=4\pi\left(\frac{m}{2\pi k_BT}\right)^{3/2}v^2e^{-mv^2/(2k_BT)},
\qquad v>0.
\]
由
\[
v=\sqrt{\frac{2\varepsilon}{m}},
\qquad \frac{dv}{d\varepsilon}=\frac{1}{\sqrt{2m\varepsilon}},
\]
得动能密度
\[
f_{\varepsilon}(\varepsilon)=f_v(v)\left|\frac{dv}{d\varepsilon}\right|
=\boxed{\frac{2}{\sqrt\pi}\frac{\sqrt\varepsilon}{(k_BT)^{3/2}}e^{-\varepsilon/(k_BT)}},
\qquad \varepsilon>0.
\]
这是形状参数 $3/2$ 的 Gamma 分布，因此
\ansmath{E(\varepsilon)=\frac32k_BT,
\qquad
\operatorname{Var}(\varepsilon)=\frac32(k_BT)^2.}
\end{solution}
\end{question}

\begin{question}
容器壁上开有一很小的小孔，外部是真空。设容器内气体处于温度 $T$ 的平衡态。
\begin{enumerate}[label=（\arabic*）,leftmargin=2.8em,itemsep=0.2em]
\item 写出穿过小孔的分子速率分布；
\item 求逸出分子的平均速率，并与容器内分子的平均速率比较。
\end{enumerate}
\begin{solution}

\begin{enumerate}[label=（\arabic*）,leftmargin=2.8em,itemsep=0.2em]
\item 逸出几率与法向速度成正比，因此逸出分子的速率分布与 $v f_v(v)$ 成正比。归一化后得
\ansmath{f_{\rm eff}(v)=2\left(\frac{m}{2k_BT}\right)^2 v^3 e^{-mv^2/(2k_BT)},
\qquad v>0.}

\item 平均速率为
\[
\langle v\rangle_{\rm eff}=\int_0^{\infty}v f_{\rm eff}(v)\,dv
=\boxed{\frac{3\sqrt\pi}{4}\sqrt{\frac{2k_BT}{m}}}.
\]
而容器内普通 Maxwell 分布的平均速率为
\[
\langle v\rangle=\sqrt{\frac{8k_BT}{\pi m}}.
\]
二者之比为
\[
\frac{\langle v\rangle_{\rm eff}}{\langle v\rangle}=\frac{3\pi}{8}>1.
\]
因此逸出分子平均更快。
\end{enumerate}
\end{solution}
\end{question}

\begin{question}
设理想气体摩尔热容 $C_V$ 为常数。推导准静态绝热过程满足的关系式
\[
pV^{\gamma}=\text{常数},
\]
并求刚性双原子理想气体从体积 $V_0$ 压缩到 $V_0/4$ 时温度比 $T_f/T_i$。
\begin{solution}

绝热过程中 $\delta Q=0$，故
\[
dU+p\,dV=0.
\]
对理想气体，$dU=nC_VdT$ 且 $pV=nRT$，因此
\[
nC_VdT+\frac{nRT}{V}dV=0.
\]
整理得
\[
\frac{dT}{T}+\frac{R}{C_V}\frac{dV}{V}=0.
\]
积分可得
\[
TV^{R/C_V}=\text{常数}.
\]
利用 $\gamma=1+R/C_V$，于是
\ansmath{pV^{\gamma}=\text{常数}.}
对刚性双原子理想气体，$f=5$，故
\[
\gamma=\frac75.
\]
由 $TV^{\gamma-1}=\text{常数}$，得
\[
\frac{T_f}{T_i}=\left(\frac{V_i}{V_f}\right)^{\gamma-1}=4^{2/5}.
\]
故
\ansmath{\frac{T_f}{T_i}=4^{2/5}.}
\end{solution}
\end{question}

\begin{question}
设大气温度恒定为 $T$，但重力加速度随高度变化为
\[
g(z)=g_0\left(\frac{R}{R+z}\right)^2,
\]
其中 $R$ 为地球半径常量。求压强分布 $p(z)$。
\begin{solution}

静力平衡给出
\[
\frac{dp}{dz}=-\rho g(z).
\]
理想气体状态方程写成
\[
\rho=\frac{\mu p}{RT},
\]
其中 $\mu$ 为摩尔质量。故
\[
\frac{dp}{p}=-\frac{\mu g_0}{RT}\left(\frac{R}{R+z}\right)^2dz.
\]
从 $0$ 积分到 $z$：
\[
\ln\frac{p(z)}{p_0}=-\frac{\mu g_0}{RT}\int_0^z\frac{R^2}{(R+\xi)^2}d\xi.
\]
而
\[
\int_0^z\frac{R^2}{(R+\xi)^2}d\xi=\frac{Rz}{R+z}.
\]
故
\ansmath{p(z)=p_0\exp\left[-\frac{\mu g_0R}{RT}\frac{z}{R+z}\right].}
\end{solution}
\end{question}


